\documentclass[../main.tex]{subfiles}

\begin{document}
\section{\textbf{Basic Concept}}
\makebox[\textwidth][s]{Author: Xianchao \hfill Date: \today}

\subsection{Schrodinger Equation}
Here we don't derive the basic Schrodinger equation, only use it as the basic axiom. Assuming the particle can be represented as a complex function, which is $\Psi(\bm{r}, t)$. The particle's position is described probabilistically, and the specific probability density at a given point in space $\bm{r}_0$ is $\rho(\bm{r}_0, t) = \Psi^*(\bm{r}_0, t) \Psi(\bm{r}_0, t) = |\Psi(\bm{r}_0, t)|^2$. The evolution of the particle can be described by time dependent Schrodinger equation:
\begin{equation}
    i\hbar \pdv{\Psi}{t} = \hat{H} \Psi
\end{equation}
The $\hat{H}$ operator is $\frac{\hat{P}^2}{2m} + \hat{V}$. Assuming the wavefunction can be written in the form as $\Psi = \textrm{e}^{-\frac{i E t}{\hbar}}\psi$, where $E$ is the energy. $\psi$ and $\hat{V}$ are independent on time. This form indicate that, the state only rotates in the complex space because the time only contribute to a change on the phase factor. Then the time dependent Schrodinger equation can be written as the time independent form:
\begin{align}
    i\hbar \pdv{(\textrm{e}^{-\frac{i E t}{\hbar}}\psi)}{t} &= \textrm{e}^{-\frac{i E t}{\hbar}} \hat{H}\psi\nonumber\\
    \textrm{e}^{-\frac{i E t}{\hbar}} E\psi &= \hat{H}\psi\nonumber\\
    E\psi &= \hat{H}\psi
\end{align}
This is the time independent Schrodinger equation.

\section{Probability Current}
Since the position of the particle can be described by probability distribution, then we can define a quantity, probability current $\bm{J}$ to denote that.
\end{document}
